\section{引言：为什么 Deep Agent 的评测如此困难？}

随着 LLM Agent 从简单的``问答式''应用演变为能够长时间自主运行、维护复杂状态的``Deep Agent''，评测方法也必须相应升级。LangChain 团队在过去一个月内发布了四个基于 Deep Agent harness 构建的应用：

\begin{itemize}
  \item \textbf{DeepAgents CLI}：编程 Agent
  \item \textbf{LangSmith Assist}：LangSmith 平台内的辅助 Agent
  \item \textbf{Personal Email Assistant}：能从交互中学习用户偏好的邮件助手
  \item \textbf{Agent Builder}：由 meta deep agent 驱动的无代码 Agent 构建平台
\end{itemize}

在为这些应用编写评测用例的过程中，LangChain 团队总结出五大核心模式，涵盖了从\textbf{单步验证}到\textbf{多轮对话模拟}、从\textbf{轨迹检查}到\textbf{环境隔离}的完整评测体系。

\begin{importantbox}{Deep Agent 与传统 LLM 应用的根本区别}
传统 LLM 评测假设每个测试用例可以用\textbf{相同的评测器}打分；而 Deep Agent 由于具备状态管理、工具调用、多步推理等能力，每个测试用例可能需要\textbf{完全不同的断言逻辑（assertion logic）}。这是本文所有讨论的出发点。
\end{importantbox}

\subsection{术语定义}

在进入具体评测模式之前，需要先统一本文使用的关键术语。LangChain 团队从两个维度对评测进行分类：\textbf{Agent 的运行方式}和\textbf{可测试的对象}。

\begin{table}[H]
\centering
\caption{Agent 运行方式分类}
\begin{tabular}{lp{10cm}}
\toprule
\textbf{运行方式} & \textbf{说明} \\
\midrule
Single Step（单步） & 限制 Agent 核心循环只执行一步，确定下一个要执行的动作 \\
Full Turn（完整轮次） & 让 Agent 在单个输入上完整运行，可能包含多次工具调用迭代 \\
Multiple Turns（多轮） & 让 Agent 多次完整运行，模拟用户与 Agent 之间的多轮对话 \\
\bottomrule
\end{tabular}
\end{table}

\begin{table}[H]
\centering
\caption{可测试对象分类}
\begin{tabular}{lp{10cm}}
\toprule
\textbf{测试对象} & \textbf{说明} \\
\midrule
Trajectory（轨迹） & Agent 调用的工具序列及其参数 \\
Final Response（最终响应） & Agent 返回给用户的最终回复 \\
Other State（其他状态） & Agent 运行过程中产生的文件、artifacts 等副产物 \\
\bottomrule
\end{tabular}
\end{table}

\begin{knowledgebox}{为什么需要区分运行方式和测试对象？}
这两个维度是正交的。同一个测试用例可能在 Full Turn 模式下运行，但只检查 Trajectory；也可能在 Single Step 模式下运行，同时检查 Final Response 和 Other State。理解这种正交关系有助于设计出更精准的测试矩阵。
\end{knowledgebox}

\subsection{本章小结}

Deep Agent 评测的核心挑战在于\textbf{状态性（statefulness）}和\textbf{长程性（long-horizon）}。传统的 input-output 评测范式无法覆盖 Agent 的工具调用轨迹、中间状态和副产物。LangChain 提出的五大模式从不同粒度和视角切入，构成了一套相对完整的评测框架。


